\chapter{uafR Core GC-MS Workflow}

\WeekHeader{4}{uafR Core GC-MS Workflow}{Convert GC-MS-style hit tables into a form that supports targeted extraction.}{Processed example workflow using \texttt{spreadOut()}, \texttt{mzExacto()}, and \texttt{exactoThese()}.}

\section{Why This Matters}
Before enrichment, students need to understand the chemical identity table coming from a mass spectrometry workflow. Tentative compound names, retention time, base peak, match factor, file name, and peak area must be reshaped carefully so chemical evidence can be compared across samples.

\section{Required Input Columns}
The core workflow expects a GC-MS-style input table with these columns:
\begin{itemize}
  \item \texttt{Component.RT}
  \item \texttt{Component.Area}
  \item \texttt{Base.Peak.MZ}
  \item \texttt{File.Name}
  \item \texttt{Compound.Name}
  \item \texttt{Match.Factor}
\end{itemize}

\section{Guided Lab}
Use the package example data first. Students should not begin with their own messy file until they can run and explain the example. The support script uses bundled exact-match output by default so the classroom workflow does not depend on live services.

\begin{rcode}
library(uafR)
source("training/scripts/04_core_workflow.R")

core <- run_core_workflow(live_lookup = FALSE)

names(core$spread)
head(core$exact)
\end{rcode}

When services are available and live lookup behavior needs to be tested, run:

\begin{rcode}
live_core <- run_core_workflow(live_lookup = TRUE)
head(live_core$exact)
\end{rcode}

Then subset an existing categorated object:

\begin{rcode}
data("standard_categorated", package = "uafR")
subset_reactives <- exactoThese(
  standard_categorated,
  subsetBy = "Database",
  subsetArgs = "reactives"
)
subset_reactives
\end{rcode}

\section{Interpretation Practice}
Students should describe:
\begin{itemize}
  \item Which compounds were retained and why.
  \item Whether match factors are high enough for the intended claim.
  \item Whether retention time or base peak patterns suggest duplicated or ambiguous hits.
  \item What should be considered tentative until confirmed by standards.
\end{itemize}

\section{Independent Extension}
Students adapt the workflow to their own data or a program-provided example. They must document every cleaning step and preserve the original input file unchanged.

\section{Deliverables}
\begin{tabularx}{\textwidth}{p{0.25\textwidth}YY}
\DeliverableTableHeader
Core workflow script & Runs from project root and produces spread and exact outputs. & R script\\
Input column audit & Confirms required columns, missing values, duplicates, and unusual match factors. & CSV or log\\
Interpretation note & Explains what can and cannot be concluded from the processed hits. & Research log\\
\end{tabularx}

\section{Quality Checklist}
\begin{checklistbox}{Before Week 5}
\begin{itemize}
  \item The script preserves raw data and writes processed outputs separately.
  \item The student can explain why tentative library matches are not structural confirmation.
  \item Duplicate or suspicious hit keys are documented.
  \item The workflow uses clear object names and comments only where they clarify decisions.
\end{itemize}
\end{checklistbox}
